% Generated draft: requires mathematical and source review.
\documentclass[11pt]{article}
\usepackage[margin=25mm]{geometry}
\usepackage{fontspec}
\setmainfont{Latin Modern Roman}
\usepackage{amsmath,amssymb,mathtools,unicode-math}
\setmathfont{Latin Modern Math}
\usepackage{enumitem,needspace,xurl,hyperref}
\setlength{\parindent}{0pt}
\setlength{\parskip}{5pt}
\setlength{\emergencystretch}{3em}
\newcommand{\N}{\mathbb{N}}
\newcommand{\Z}{\mathbb{Z}}
\newcommand{\Q}{\mathbb{Q}}
\newcommand{\R}{\mathbb{R}}
\newcommand{\C}{\mathbb{C}}
\newcommand{\F}{\mathbb{F}}
\newcommand{\A}{\mathbb{A}}
\newcommand{\PP}{\mathbb{P}}
\newcommand{\E}{\mathbb{E}}
\newcommand{\eps}{\varepsilon}
\newcommand{\dd}{\,\mathrm{d}}
\begin{document}
\section*{UnsolvedMath ID: 701772}
\textbf{Source problem code:} catalogue entry
\subsection{Definitions and mathematical statement}
For a finite group \(G\), its exponent \(\exp(G)\) is the least positive integer \(q\) such that \(g^q=1\) for every \(g\in G\). Its Schur multiplier is the finite abelian group
\[
M(G)=H_2(G,\Z),
\]
where group homology is taken with trivial coefficients. The exponent of the trivial group is taken to be \(1\).

The question proposed by M. R. Vaughan-Lee is whether
\[
\exists c\in\Z_{>0}\;\forall q\in\Z_{>0}\;\forall\text{ finite groups }G,
\qquad
\exp(G)=q\ \Longrightarrow\ \exp(M(G))\mid q^c.
\]
The integer \(c\) must be independent of both \(q\) and \(G\). Here the source phrase ``for any \(q\)'' is interpreted as quantification over positive integers, as appropriate for finite group exponents. The supplied registry labels this question open; that label is not an independently verified assessment.
\subsection{Short English statement}
Does one universal positive integer \(c\) ensure that the exponent of the Schur multiplier of every finite group of exponent \(q\) divides \(q^c\)? The bound must hold uniformly for all positive integers \(q\) and all such finite groups.
\subsection{Sources}
\begin{itemize}
\item Supplied local UnsolvedMath registry record, permanent ID \(701772\), exact source problem code \texttt{catalogue entry}, titled ``Kourovka Notebook Problem 20.114.'' The record labels the problem open and attributes it to M. R. Vaughan-Lee, with proposed year 2022.
\item The supplied background cites \emph{The Kourovka Notebook (Unsolved Problems in Group Theory)}, 20th issue (2022), and the 21st issue edition (2026), identified as \texttt{arXiv:1401.0300v48}; original problem number 20.114. The collection description names E. I. Khukhro and V. D. Mazurov as editors. These bibliographic details are supplied attribution and have not been independently verified.
\end{itemize}
\end{document}
